\section{The exponential distribution}
\label{sec:exp-dist}

We say that $X$ has the exponential distribution with
parameter $\lambda$, and denote
$X\sim \operatorname{Exp}(\lambda)$, if $X$ has
density function

\[
f_X(x) = \lambda e^{-\lambda x} \qquad (x\ge 0)
\]

\noindent and the associated c.d.f.

\[
F_X(x) = \begin{cases} 0 & \textrm{if }x<0, \\
1-e^{-\lambda x} & \textrm{if }x\ge 0.
\end{cases}
\]

\noindent $\textbf{Properties:}$

\begin{enumerate}[label=\arabic*.]
\item $\lambda$ has the role of an (inverse)
$\textbf{scale parameter}$, in the sense that for
$c>0$ we have that
$c \operatorname{Exp}(\lambda) \stackrel{d}{=}
\operatorname{Exp}(\lambda / a)$;
i.e., scaling an exponential r.v. by a factor $c$
gives a new exponential r.v. where the scale parameter
is divided by $c$.

\item The exponential distribution satisfies the
$\textbf{lack of memory property}$. More precisely,
if $X\sim \operatorname{Exp}(\lambda)$ then

\[
\mathrm{P}(X > t+s \,|\, X>t) = \mathrm{P}(X>s) \qquad (t,s>0).
\]

\noindent Furthermore, it is not hard to show that the
exponential distribution is the unique distribution on
$[0,\infty)$ satisfying this property.

\item $\mathbf{E}(X) = \lambda$,
$\mathbf{V}(X) = \lambda^2$.

\item The exponential distribution can be thought of as a
scaling limit of geometric random variables, when the
geometric distribution is interpreted as measuring
\emph{time} rather than the number of experiments, and time
is scaled so that the i.i.d. Bernoulli experiments are
performed more and more frequently, but are becoming
less and less probable to succeed, in such a way that
the mean number of successful experiments per unit of
time remains constant. More precisely, if $\lambda>0$
is fixed, $X\sim \operatorname{Exp}(\lambda)$, and
for each $n$ (larger than $\lambda$) we let $W_n$
denote a random variable with distribution
$\operatorname{Geom}(\lambda/n)$, then we have

\[
\mathrm{P}\left(\frac{1}{n} W_n > t\right)  \xrightarrow[n\to\infty]{} \mathrm{P}(X > t) \qquad (t>0).
\]

\item If $X\sim \operatorname{Exp}(\lambda)$ and
$Y \sim \operatorname{Exp}(\mu)$ are independent
r.v.s then
$\min(X,Y) \sim \operatorname{Exp}(\lambda+\mu)$.

\item If $X_1,X_2,\ldots$ are i.i.d. $\operatorname{Exp}(1)$
random variables, and we define the cumulative sums
$S_0=0$, $S_n = \sum_{k=1}^n X_k$, then for each
$\lambda>0$, the random variable

\[
N(\lambda) = \max\{ n\ge 0\,:\, S_n \le \lambda \},
\]

\noindent then $N(\lambda) \sim \operatorname{Poi}(\lambda)$.
\end{enumerate}

\noindent $\textbf{Note.}$ One can consider $N(\lambda)$ not
just for a single value of $\lambda$ but the entire
family $N(t)$, where $t>0$ is a parameter denoting
time. When considered as such a family, $(N(t))_{t>0}$
is called a $\textbf{Poisson process}$.
